% SIAM Article Template
% Preserve modern LaTeX reference commands before the legacy SIAM class loads.
% Load this file BEFORE \\documentclass; remove when upgrading the SIAM class.
\let\SIAMkernelLabel\label
\let\SIAMkernelRefstepcounter\refstepcounter
\AddToHook{package/hyperref/before}{%
  \let\label\SIAMkernelLabel
  \let\refstepcounter\SIAMkernelRefstepcounter
}

\documentclass[review,hidelinks,onefignum,onetabnum]{siamart220329}

% Information that is shared between the article and the supplement
% (title and author information, macros, packages, etc.) goes into
% ex_shared.tex. If there is no supplement, this file can be included
% directly.

% !TeX root = ex_article.tex
% Shared packages, notation, and manuscript information.


% Packages and macros go here
\usepackage{lipsum}
\usepackage{amsfonts}
\usepackage{graphicx}
\usepackage{epstopdf}
\usepackage{algorithmic}
\ifpdf
  \DeclareGraphicsExtensions{.eps,.pdf,.png,.jpg}
\else
  \DeclareGraphicsExtensions{.eps}
\fi

% Add a serial/Oxford comma by default.
\newcommand{\creflastconjunction}{, and~}

% Used for creating new theorem and remark environments
\newsiamremark{remark}{Remark}
\newsiamremark{hypothesis}{Hypothesis}
\crefname{hypothesis}{Hypothesis}{Hypotheses}
\newsiamthm{claim}{Claim}

% Short running headers for the article and supplement.
\headers{SDIRK mr-ccSAV Schemes}
        {H. Wang and X. Wang}

% Title. If the supplement option is on, then "Supplementary Material"
% is automatically inserted before the title.
\newcommand{\papertitle}{%
  Long-Time Stability of SDIRK mr-ccSAV Schemes
  for the Forced Two-Dimensional Navier--Stokes Equations
  with No-Slip Boundary Conditions}
\title{\papertitle}

% Authors: full names plus addresses.
\author{Haifeng Wang\thanks{%
  School of Mathematical Sciences, Eastern Institute of Technology,
  Ningbo, Zhejiang 315200, China
  (\email{hfwang@eitech.edu.cn}, \email{wxm@eitech.edu.cn}).}
  \and Xiaoming Wang\footnotemark[1]}

% Keep PDF metadata consistent with the displayed manuscript information.
\ifpdf
\hypersetup{
  pdftitle={\siampretitle\papertitle},
  pdfauthor={Haifeng Wang and Xiaoming Wang}
}
\fi

\usepackage{amsopn}
\DeclareMathOperator{\diag}{diag}

% \RequirePackage[top=2cm,bottom=2cm,left=3cm,right=3cm,includehead,includefoot]{geometry}


%%% Local Variables: 
%%% mode:latex
%%% TeX-master: "ex_article"
%%% End: 


% Optional PDF information
\ifpdf
\hypersetup{
  pdftitle={An Example Article},
  pdfauthor={D. Doe, P. T. Frank, and J. E. Smith}
}
\fi

% The next statement enables references to information in the
% supplement. See the xr-hyperref package for details.

\externaldocument[][nocite]{ex_supplement}

% FundRef data to be entered by SIAM
%<funding-group specific-use="FundRef">
%<award-group>
%<funding-source>
%<named-content content-type="funder-name"> 
%</named-content> 
%<named-content content-type="funder-identifier"> 
%</named-content>
%</funding-source>
%<award-id> </award-id>
%</award-group>
%</funding-group>

\begin{document}

\maketitle

% REQUIRED
\begin{abstract}
This is an example SIAM \LaTeX\ article. This can be used as a
template for new articles.  Abstracts must be able to stand alone
and so cannot contain citations to the paper's references,
equations, etc.  An abstract must consist of a single paragraph and
be concise. Because of online formatting, abstracts must appear as
plain as possible. Any equations should be inline.
\end{abstract}

% REQUIRED
\begin{keywords}
example, \LaTeX
\end{keywords}

% REQUIRED
\begin{MSCcodes}
68Q25, 68R10, 68U05
\end{MSCcodes}

\section{Introduction}
The incompressible Navier-Stokes equations (NSE) are a fundamental model of incompressible fluid flow and provide a canonical testbed for numerical methods designed for long-time simulations. Let $\Omega\subset\mathbb R^2$ be a bounded, connected domain with a $C^2$ boundary. We focus in this paper the forced incompressible NSE in dimensionless, subject to the no-slip boundary condition:
\begin{equation}\label{eq:nse}
\begin{aligned}
 u_t-\nu\Delta u+(u\cdot\nabla)u+\nabla p&=f
 &&\text{in }\Omega\times(0,\infty),\\
 \nabla\cdot u&=0
 &&\text{in }\Omega\times(0,\infty),\\
 u&=0 &&\text{on }\partial\Omega\times(0,\infty),\\
 u(\cdot,0)&=u_0 &&\text{in }\Omega.
\end{aligned}
\end{equation}
Here, $u$ and $p$ denote the velocity and pressure, respectively, $\nu>0$ is the kinematic viscosity, $f$ is the external body force, and $u_0$ is the initial velocity. The pressure is normalized by requiring $\int_\Omega p(x,t)\,dx=0$ for each $t>0$.

\section{Preliminaries and the mean reverting ccSAV reformulation of the NSE}

\subsection{Function space and notation}\label{sec:functional-setting}
For the mathematical setting of the problem, we consider the following space
\begin{equation}
  \begin{aligned}
 &V=\{v\in H_0^1(\Omega)^2: \nabla\cdot v=0\}, \\
 &H=\{v \in L^2(\Omega)^2: \nabla \cdot v=0, v\cdot n = 0 \text{ on } \partial\Omega\}
  \end{aligned}
\end{equation}
We write $(\cdot,\cdot)$ and $|\cdot|$ for the $L^2$ inner product and
norm, and equip $V$ with the norm $\|v\|=|\nabla v|$.
The pairing between $V'$ and $V$ is denoted by $\langle\cdot,\cdot\rangle$.
Thus $V\subset H\equiv H'\subset V'$. The full homogeneous velocity
trace is imposed in $V$; membership in $H$ alone does not impose a
zero tangential trace. We also use
$L_0^2(\Omega)=\{q\in L^2(\Omega):\int_\Omega q\,dx=0\}$.

Let $P:L^2(\Omega)^2\to H$ be the orthogonal Leray projection.
The Stokes operator and its domain are
\[
 A=-P\Delta,\qquad D(A)=H^2(\Omega)^2\cap V.
\]
Its extension from $V$ to $V'$ is defined by
$\langle Au,v\rangle=(\nabla u,\nabla v)$.
Writing $\lambda_1>0$ for its first eigenvalue, we have
\begin{equation}\label{eq:stokes-estimates}
 \lambda_1|v|^2\le\|v\|^2\quad(v\in V),\qquad
 \|v\|_{H^2}\le C|Av|\quad(v\in D(A)).
\end{equation}
Here and below, $C$ denotes a positive constant depending only on the
fixed domain unless further dependencies are indicated; its value may
change between occurrences.

For $u,v,w\in V$, define
\[
 b(u,v,w)=\int_\Omega (u\cdot\nabla)v\cdot w\,dx,
 \qquad \langle B(u,v),w\rangle=b(u,v,w).
\]
Integration by parts gives
\begin{equation}\label{eq:skew-symmetry}
 b(u,v,w)=-b(u,w,v),\qquad b(u,v,v)=0.
\end{equation}
The two-dimensional interpolation inequalities and Stokes regularity
also give
\begin{align}
 |b(u,v,w)|&\le C|u|^{1/2}\|u\|^{1/2}\|v\|
                      |w|^{1/2}\|w\|^{1/2},
 &&u,v,w\in V,\label{eq:trilinear-v}\\
 |b(u,v,w)|&\le C|u|^{1/2}|Au|^{1/2}\|v\||w|,
 &&u\in D(A),\ v\in V,\ w\in H.\label{eq:trilinear-a}
\end{align}
In the second estimate, the integral defining $b$ extends to the stated
spaces. In particular, $B:V\times V\to V'$ is well defined, and
$B(u,v)=P((u\cdot\nabla)v)$ whenever the convective term belongs to
$L^2(\Omega)^2$. This does not imply that $B(u,v)$ belongs to $V$.
These are the standard no-slip functional setting and estimates; see,
for example, \cite[section~1]{ToneWirosoetisno2006}.

Throughout this section, assume
\begin{equation}\label{eq:continuous-data}
 u_0\in V,\qquad f\in L^\infty(0,\infty;L^2(\Omega)^2),\qquad
 F=Pf,\qquad F_* =\operatorname*{ess\,sup}_{t\ge0}|F(t)|.
\end{equation}
Applying $P$ to \cref{eq:nse} gives
\begin{equation}\label{eq:projected-nse}
 u_t+\nu Au+B(u,u)=F,\qquad u(0)=u_0.
\end{equation}
The unique two-dimensional strong solution belongs to
$C([0,T];V)\cap L^2(0,T;D(A))$, with $u_t\in L^2(0,T;H)$,
for every finite $T>0$, and satisfies the kinetic energy identity
\begin{equation}\label{eq:nse-energy}
 \frac12\frac{d}{dt}|u|^2+\nu\|u\|^2=(F,u).
\end{equation}
We recall the classical uniform-in-time bounds below; see
\cite[section~1, equations~(1.6)--(1.8)]{ToneWirosoetisno2006}.
The accompanying time-window estimate follows from the energy balance.
We state these results without proof, keeping their dependence on the
initial data explicit.

\begin{proposition}[Uniform bounds for the NSE]\label{prop:nse-bounds}
Under \cref{eq:continuous-data}, the solution of
\cref{eq:projected-nse} satisfies, for $t\ge0$,
\begin{equation}\label{eq:nse-l2-bound}
 |u(t)|^2\le e^{-\nu\lambda_1t}|u_0|^2
 +\frac{F_*^2}{\nu^2\lambda_1^2}
       (1-e^{-\nu\lambda_1t}).
\end{equation}
For every fixed $\ell>0$ and every $t\ge0$,
\begin{equation}\label{eq:nse-h1-window}
 \nu\int_t^{t+\ell}\|u(s)\|^2\,ds
 \le |u(t)|^2+\frac{\ell F_*^2}{\nu\lambda_1}.
\end{equation}
Moreover, there exists $K_1=K_1(\nu,\Omega,F_*,\|u_0\|)$ such that
\begin{equation}\label{eq:nse-h1-bound}
 \sup_{t\ge0}\|u(t)\|^2\le K_1.
\end{equation}
There are radii $R_0,R_1>0$, depending only on $\nu,\Omega,F_*$, and
an entry time $T_*=T_*(\nu,\Omega,F_*,|u_0|)$ such that
\begin{equation}\label{eq:nse-absorbing}
 |u(t)|^2\le R_0^2,\qquad \|u(t)\|^2\le R_1^2,
 \qquad t\ge T_*.
\end{equation}
The entry time can be chosen uniformly for initial data in a bounded
subset of $H$ that is contained in $V$.
\end{proposition}


\subsection{The mean-reverting ccSAV reformulation}
\label{sec:mr-ccsav}

The cancellation in \cref{eq:skew-symmetry} underlies the continuous
energy estimate. In a time discretization with an explicitly evaluated
convective term, its pairing with the current stage velocity need not
vanish. To retain an algebraic cancellation for such pairings, introduce
a scalar variable $r$ and a mean-reversion parameter $\gamma>0$, and set
\[
 G_u(r)=1-r^2,\qquad G_r(r)=1+r.
\]
The mean-reverting concurrently corrected scalar auxiliary variable
(mr-ccSAV) system is
\begin{equation}\label{eq:mr-ccsav}
\begin{aligned}
 u_t+\nu Au+(1-r^2)B(u,u)&=F,\\
 r_t+\gamma r+(1+r)\langle B(u,u),u\rangle&=0,\\
 u(0)=u_0,\qquad r(0)&=0.
\end{aligned}
\end{equation}
The force remains in the velocity equation and represents physical
energy input. The scalar correction involves only the convective term.

\begin{proposition}[Equivalence and modified energy identity]
\label{prop:mr-equivalence}
Under \cref{eq:continuous-data}, system \cref{eq:mr-ccsav} is equivalent
to \cref{eq:projected-nse} in the strong-solution class, with $r\equiv0$.
For
\begin{equation}\label{eq:modified-energy}
 \mathcal E(u,r)=\frac12|u|^2+\frac12(r-1)^2,
\end{equation}
its solution satisfies
\begin{equation}\label{eq:mr-energy}
 \frac{d}{dt}\mathcal E(u,r)+\nu\|u\|^2+\gamma r(r-1)=(F,u),
\end{equation}
or, equivalently,
\begin{equation}\label{eq:mr-energy-shifted}
 \frac{d}{dt}\mathcal E(u,r)+\nu\|u\|^2
 +\frac\gamma2\bigl(r^2+(r-1)^2\bigr)
 =(F,u)+\frac\gamma2.
\end{equation}
\end{proposition}

\begin{proof}
By \cref{eq:skew-symmetry}, the scalar equation reduces to
$r_t+\gamma r=0$. Its initial condition implies $r(t)=0$, and the
velocity equation becomes \cref{eq:projected-nse}. Conversely, its
strong solution together with $r=0$ solves \cref{eq:mr-ccsav}.

For the energy identity, pair the velocity equation with $u$ and
multiply the scalar equation by $r-1$. Their convective contributions
cancel because
\begin{equation}\label{eq:ccsav-cancellation}
 G_u(r)+(r-1)G_r(r)=(1-r^2)+(r-1)(1+r)=0.
\end{equation}
Adding the two equations yields \cref{eq:mr-energy}. The identity
$2r(r-1)=r^2+(r-1)^2-1$ gives \cref{eq:mr-energy-shifted}.
\end{proof}

The equilibrium of the auxiliary variable is $r=0$, although its
contribution to the modified energy is centered at $r=1$. At $r=0$,
\cref{eq:mr-energy} reduces exactly to the physical energy balance
\cref{eq:nse-energy}. More generally, the homogeneous scalar equation
has the decay $r(t)=r(0)e^{-\gamma t}$, which explains the mean-reverting
term. The velocity correction is quadratic in the deviation from this
equilibrium: $G_u(r)-1=-r^2$.

The algebraic identity \cref{eq:ccsav-cancellation} remains valid if
$B(u,u)$ in both equations is replaced by a common explicitly evaluated
convective term, provided that the same scalar and velocity are used
in the pairing. This property motivates the concurrent stage
correction in the SDIRK2 scheme below. It supplies the cancellation
needed for the discrete energy estimate; the bounds on the numerical
auxiliary variable and the temporal accuracy require separate analysis.

\section{The SDIRK2-mr-ccSAV scheme and its solvability}
\label{sec:alg}

In this section, we construct the SDIRK2-mr-ccSAV scheme in incremental
form. Each stage is reduced to two linear Stokes problems with the
same operator and a scalar polynomial equation. This reduction yields
existence for every positive time step and a computable sufficient
condition for uniqueness.

\subsection{The time-stepping scheme}
\label{sec:sdirk2-scheme}

Let $0=t_0<t_1<\cdots$ be a time grid, with
$\tau_n=t_{n+1}-t_n>0$. Set
\begin{equation}\label{eq:sdirk-parameters}
 \eta=1-\frac1{\sqrt2},\qquad
 \delta=1-\frac1{2\eta}=-\frac1{\sqrt2},\qquad
 (c_0,c_1,c_2)=(0,\eta,1).
\end{equation}
The coefficient matrices for the two nontrivial stages of the underlying
implicit--explicit method are
\[
 \mathsf A^{\rm RK}=\begin{pmatrix}\eta&0\\1-\eta&\eta\end{pmatrix},
 \qquad
 \widehat{\mathsf A}^{\rm RK}
 =\begin{pmatrix}\eta&0\\\delta&1-\delta\end{pmatrix}.
\]
The implicit columns correspond to stages $1,2$, whereas the explicit
columns correspond to stages $0,1$. Taking differences of successive
rows gives the incremental coefficients
\begin{equation}\label{eq:sdirk-incremental}
 \mathsf A=(a_{ij})
 =\begin{pmatrix}\eta&0\\1-2\eta&\eta\end{pmatrix},\qquad
 \widehat{\mathsf A}=(\widehat a_{ij})
 =\begin{pmatrix}\eta&0\\-1&1-\delta\end{pmatrix},
\end{equation}
where $\delta-\eta=-1$. The concurrent correction below uses the
current scalar in each incremental stage equation.

Initialize $u^0=u_0$ and $r^0=0$. Given $(u^n,r^n)$, set
$(u_{n,0},r_{n,0})=(u^n,r^n)$ and write
$B_{n,j}=B(u_{n,j},u_{n,j})$. Let $F_{n,j}$ denote the prescribed stage
force. For a continuous force, take $F_{n,j}=F(t_n+c_j\tau_n)$,
$j=0,1$. If $F$ is given only as an $L^\infty$ equivalence class,
the stage data must be specified separately, for example by fixing a
bounded representative; solvability below concerns the actual data
$F_{n,j}$ used in the scheme.

For $i=1,2$, find $(u_{n,i},r_{n,i})\in V\times\mathbb R$ such that
\begin{equation}\label{eq:sdirk-mr-scheme}
\begin{aligned}
 u_{n,i}-u_{n,i-1}
 +\nu\tau_n\sum_{j=1}^i a_{ij}Au_{n,j}
 &=\tau_n\sum_{j=0}^{i-1}\widehat a_{ij}
       \bigl[F_{n,j}-(1-r_{n,i}^2)B_{n,j}\bigr],\\
 r_{n,i}-r_{n,i-1}
 +\gamma\tau_n\sum_{j=1}^i a_{ij}r_{n,j}
 &=-\tau_n(1+r_{n,i})\sum_{j=0}^{i-1}\widehat a_{ij}
       \langle B_{n,j},u_{n,i}\rangle.
\end{aligned}
\end{equation}
The first equation is understood in $V'$. The step is completed by
\begin{equation}\label{eq:sdirk-update}
 (u^{n+1},r^{n+1})=(u_{n,2},r_{n,2}).
\end{equation}
Both correction factors in a given stage use $r_{n,i}$, while the
convective terms depend only on already computed velocities. Consequently,
\cref{eq:ccsav-cancellation} applies to the common pairing
$\sum_{j<i}\widehat a_{ij}\langle B_{n,j},u_{n,i}\rangle$.
The corrected incremental scheme is not obtained simply by applying
the displayed Runge--Kutta matrices to the continuous extended system.

\subsection{Reduction to a scalar equation}
\label{sec:scalar-reduction}

Fix a stage $i$ and regard all preceding stages as known. Define
\begin{equation}\label{eq:stage-data}
\begin{aligned}
 \widetilde B_{n,i}
 &=\sum_{j=0}^{i-1}\widehat a_{ij}B_{n,j},\\
 R^u_{n,i}
 &=u_{n,i-1}-\nu\tau_n\sum_{j=1}^{i-1}a_{ij}Au_{n,j}
       +\tau_n\sum_{j=0}^{i-1}\widehat a_{ij}F_{n,j},\\
 R^r_{n,i}
 &=r_{n,i-1}-\gamma\tau_n\sum_{j=1}^{i-1}a_{ij}r_{n,j}.
\end{aligned}
\end{equation}
Empty sums are zero. In particular, the first-stage data are
\[
 \widetilde B_{n,1}=\eta B_{n,0},\qquad
 R^u_{n,1}=u^n+\eta\tau_nF_{n,0},\qquad R^r_{n,1}=r^n,
\]
and the second-stage data are
\[
\begin{aligned}
 \widetilde B_{n,2}&=-B_{n,0}+(1-\delta)B_{n,1},\\
 R^u_{n,2}&=u_{n,1}-\nu\tau_n(1-2\eta)Au_{n,1}
                   +\tau_n[-F_{n,0}+(1-\delta)F_{n,1}],\\
 R^r_{n,2}&=[1-\gamma\tau_n(1-2\eta)]r_{n,1}.
\end{aligned}
\]
Thus \cref{eq:sdirk-mr-scheme} takes the form
\begin{equation}\label{eq:stage-resolvent}
\begin{aligned}
 L_nu_{n,i}&=R^u_{n,i}-\tau_n(1-r_{n,i}^2)\widetilde B_{n,i},\\
 (1+\gamma\eta\tau_n)r_{n,i}
 &=R^r_{n,i}-\tau_n(1+r_{n,i})
                     \langle\widetilde B_{n,i},u_{n,i}\rangle,
 \qquad L_n=I+\nu\eta\tau_n A.
\end{aligned}
\end{equation}
Here $I:V\to V'$ is the embedding induced by the $L^2$ inner product.

Solve the two linear problems
\begin{equation}\label{eq:linear-stage-solves}
 L_nv_{n,i}=R^u_{n,i},\qquad
 L_nw_{n,i}=\widetilde B_{n,i},
\end{equation}
and set
\begin{equation}\label{eq:scalar-coefficients}
 \alpha_{n,i}=\langle\widetilde B_{n,i},v_{n,i}\rangle,
 \qquad
 \beta_{n,i}=\langle\widetilde B_{n,i},w_{n,i}\rangle.
\end{equation}
For any prescribed scalar $r$, the velocity equation then has the solution
\begin{equation}\label{eq:velocity-reconstruction}
 u_{n,i}(r)=v_{n,i}-\tau_n(1-r^2)w_{n,i}.
\end{equation}
Substitution into the scalar equation gives
\begin{equation}\label{eq:stage-polynomial}
\begin{aligned}
 \mathcal F_{n,i}(r)
 &=(1+\gamma\eta\tau_n)r-R^r_{n,i}
       +\tau_n(1+r)\alpha_{n,i}\\
 &\quad+\tau_n^2\beta_{n,i}(r^3+r^2-r-1)=0.
\end{aligned}
\end{equation}
Every real root determines a stage solution through
\cref{eq:velocity-reconstruction}, and every stage solution arises in
this way. The operator $L_n$ is the same for both auxiliary solves and
both stages of a time step. At a fixed time step size it is also
unchanged between time steps.

\subsection{Existence and conditional uniqueness}
\label{sec:stage-solvability}

\begin{proposition}[Stage solvability]\label[proposition]{prop:stage-existence}
Let $u^n\in V$, $r^n\in\mathbb R$, and $F_{n,0},F_{n,1}\in V'$.
For every $\tau_n>0$, the two stages of \cref{eq:sdirk-mr-scheme}
can be solved successively, and each stage has at least one solution
in $V\times\mathbb R$. If, in addition, $u^n\in D(A)$ and
$F_{n,0},F_{n,1}\in H$, then every stage velocity belongs to $D(A)$.
\end{proposition}

\begin{proof}
For $z\in V$, skew-symmetry and the two-dimensional embedding into
$L^4$ give $\|B(z,z)\|_{V'}\le C\|z\|^2$. Also $Az\in V'$.
Hence, when the preceding velocities belong to $V$, both
$R^u_{n,i}$ and $\widetilde B_{n,i}$ belong to $V'$.
The bilinear form associated with $L_n$ is continuous on $V\times V$
and satisfies
\[
 \langle L_nz,z\rangle=|z|^2+\nu\eta\tau_n\|z\|^2
 \ge\nu\eta\tau_n\|z\|^2.
\]
The Lax--Milgram theorem therefore gives unique solutions
$v_{n,i},w_{n,i}\in V$ of \cref{eq:linear-stage-solves}.
Testing the second problem with $w_{n,i}$ yields
\begin{equation}\label{eq:beta-nonnegative}
 \beta_{n,i}=|w_{n,i}|^2+\nu\eta\tau_n\|w_{n,i}\|^2\ge0.
\end{equation}
If $\beta_{n,i}>0$, then $\mathcal F_{n,i}$ is a cubic polynomial
with positive leading coefficient and has at least one real root.
If $\beta_{n,i}=0$, \cref{eq:beta-nonnegative} implies $w_{n,i}=0$,
so that $\widetilde B_{n,i}=0$ and $\alpha_{n,i}=0$.
The scalar equation then reduces to a linear equation with the unique root
\[
 r_{n,i}=\frac{R^r_{n,i}}{1+\gamma\eta\tau_n}.
\]
In either case, \cref{eq:velocity-reconstruction} gives a stage solution
in $V\times\mathbb R$. Applying this argument first to $i=1$ and then
to $i=2$ proves existence of the complete step without a step size restriction.

For the regularity assertion, $z\in D(A)$ implies
$B(z,z)\in H$, since $H^2(\Omega)\subset L^\infty(\Omega)$ in two
dimensions. Thus $D(A)$ preceding velocities and $H$ stage forces give
$R^u_{n,i},\widetilde B_{n,i}\in H$. If $L_nz=g\in H$, its weak
solution satisfies $Az=(g-z)/(\nu\eta\tau_n)\in H$, hence
$z\in D(A)$. Both auxiliary velocities and their linear combination
\cref{eq:velocity-reconstruction} consequently belong to $D(A)$.
Induction over the two stages proves the assertion.
\end{proof}

\begin{proposition}[A sufficient condition for uniqueness]
\label[proposition]{prop:stage-uniqueness}
For fixed preceding stage data, the stage solution is unique if
\begin{equation}\label{eq:stage-unique-condition}
 1+\gamma\eta\tau_n+\tau_n\alpha_{n,i}
       -\frac43\tau_n^2\beta_{n,i}>0.
\end{equation}
In particular, the more conservative condition
\begin{equation}\label{eq:stage-unique-conservative}
 \tau_n|\alpha_{n,i}|+\frac43\tau_n^2\beta_{n,i}
       <1+\gamma\eta\tau_n
\end{equation}
is sufficient. If the first stage satisfies the condition and the
second stage satisfies it for the resulting first-stage data, the
complete time step has a unique solution.
\end{proposition}

\begin{proof}
Differentiating \cref{eq:stage-polynomial} and completing the square gives
\[
\begin{aligned}
 \mathcal F'_{n,i}(r)
 &=1+\gamma\eta\tau_n+\tau_n\alpha_{n,i}
   +\tau_n^2\beta_{n,i}\left[3\left(r+\tfrac13\right)^2-\tfrac43\right]\\
 &\ge1+\gamma\eta\tau_n+\tau_n\alpha_{n,i}
                         -\tfrac43\tau_n^2\beta_{n,i}.
\end{aligned}
\]
Under \cref{eq:stage-unique-condition}, the polynomial is strictly
increasing on $\mathbb R$. Existence from
\cref{prop:stage-existence} then gives exactly one real root.
The reconstruction formula determines the velocity uniquely.
Condition \cref{eq:stage-unique-conservative} implies
\cref{eq:stage-unique-condition}; applying the argument successively
proves the assertion for the full step.
\end{proof}

The uniqueness conditions are sufficient, not necessary: their failure
does not imply that a stage has multiple solutions. Without such a
condition, the scheme admits any real root of \cref{eq:stage-polynomial};
a numerical implementation must specify its root-selection rule when
several roots are present. The existence result holds for every such
choice. Likewise, propagation of $D(A)$ regularity does not by itself
give a time-uniform or step-size-independent $H^2$ bound.

\begin{algorithm}[htbp]
\caption{One step of the SDIRK2-mr-ccSAV scheme}
\label{alg:sdirk-mr}
\begin{algorithmic}[1]
\REQUIRE $(u^n,r^n)$, $\tau_n>0$, and stage forces $F_{n,0},F_{n,1}$.
\STATE Set $(u_{n,0},r_{n,0})=(u^n,r^n)$.
\FOR{$i=1,2$}
\STATE Assemble $\widetilde B_{n,i}$, $R^u_{n,i}$, and $R^r_{n,i}$
       from \cref{eq:stage-data}.
\STATE Solve the two linear Stokes problems \cref{eq:linear-stage-solves}.
\STATE Compute $\alpha_{n,i}$ and $\beta_{n,i}$ from
       \cref{eq:scalar-coefficients}.
\IF{$\beta_{n,i}=0$}
\STATE Set $r_{n,i}=R^r_{n,i}/(1+\gamma\eta\tau_n)$.
\ELSE
\STATE Solve \cref{eq:stage-polynomial} for a real root $r_{n,i}$,
       using the specified selection rule if necessary.
\ENDIF
\STATE Recover $u_{n,i}$ from \cref{eq:velocity-reconstruction}.
\ENDFOR
\STATE Set $(u^{n+1},r^{n+1})=(u_{n,2},r_{n,2})$.
\end{algorithmic}
\end{algorithm}

\Cref{alg:sdirk-mr} summarizes the construction in exact arithmetic.
In floating-point computations, the scalar residual should be checked
using the original scalar equation in \cref{eq:stage-resolvent}; a small computed
value of $\beta_{n,i}$ alone does not justify discarding all convective
terms.

\section{Long-time stability in \texorpdfstring{$L^2$ and $H^1$}{L2 and H1}}
\label{sec:long-time-stability}

In this section, we establish uniform-in-time bounds for the numerical
velocity and auxiliary variable. The basic energy estimate holds on an
arbitrary time grid and yields an $L^2$ absorbing bound together with
an $H^1$ time-window estimate. For a fixed time step, we then prove
uniform $H^1$ stability by controlling a modified higher-order energy
and applying a discrete uniform Gronwall argument. All estimates apply
to every real-root branch of the scheme and do not require stage
uniqueness.

\subsection{The discrete energy estimate and long-time \texorpdfstring{$L^2$}{L2} stability}
\label{sec:discrete-l2}

Write
\begin{equation}\label{eq:lambda-incremental}
 \lambda_I=2\eta-\tfrac12=\tfrac32-\sqrt2>0.
\end{equation}
For any two vectors $z_1,z_2$ in a real Hilbert space, the incremental
matrix in \cref{eq:sdirk-incremental} satisfies
\begin{equation}\label{eq:incremental-coercivity}
\begin{aligned}
 \sum_{i=1}^2\sum_{j=1}^i a_{ij}(z_j,z_i)
 &=\lambda_I\sum_{i=1}^2|z_i|^2
       +\frac{1-2\eta}{2}|z_1+z_2|^2\\
 &\ge\lambda_I\sum_{i=1}^2|z_i|^2.
\end{aligned}
\end{equation}
The same identity holds for real scalars. Define the nodal energy
\[
 E^n=\mathcal E(u^n,r^n)=\tfrac12\bigl(|u^n|^2+|r^n-1|^2\bigr).
\]
In this subsection assume $u^0\in V$, $r^0\in\mathbb R$, and
\begin{equation}\label{eq:discrete-force-dual}
 M_{F,V'}:=\sup_{n\ge0,\ j=0,1}\|F_{n,j}\|_{V'}<\infty,
 \qquad
 \|G\|_{V'}=\sup_{0\ne v\in V}\frac{|\langle G,v\rangle|}{\|v\|}.
\end{equation}
Allowing an arbitrary $r^0$ also covers the prescribed initialization
$r^0=0$. Fix the constants
\begin{equation}\label{eq:l2-constants}
\begin{aligned}
 C_0&=\frac{\gamma[\eta^2+(1-\eta)^2]}{2\lambda_I}
       +\frac{M_{F,V'}^2[\eta^2+(2-\delta)^2]}{2\nu\lambda_I},\\
 \kappa&=\lambda_I\min\{\nu\lambda_1,\gamma\},\qquad
 M_E=\max\{E^0,C_0/\kappa\}.
\end{aligned}
\end{equation}
These constants are independent of the time grid and the chosen roots.

\begin{lemma}[Discrete energy inequality]\label[lemma]{lem:discrete-energy}
For every $\tau_n>0$, the scheme satisfies
\begin{equation}\label{eq:discrete-energy-inequality}
\begin{aligned}
 E^{n+1}-E^n
 &+\frac12\sum_{i=1}^2
       \bigl(|u_{n,i}-u_{n,i-1}|^2+|r_{n,i}-r_{n,i-1}|^2\bigr)\\
 &+\frac{\nu\lambda_I}{2}\tau_n\sum_{i=1}^2\|u_{n,i}\|^2
  +\frac{\gamma\lambda_I}{2}\tau_n\sum_{i=1}^2|r_{n,i}-1|^2
 \le C_0\tau_n.
\end{aligned}
\end{equation}
\end{lemma}

\begin{proof}
Pair the velocity equation in \cref{eq:sdirk-mr-scheme} with
$u_{n,i}$ and multiply the scalar equation by $r_{n,i}-1$.
The convective terms cancel by \cref{eq:ccsav-cancellation}.
Summing over $i=1,2$ and using
$2(x-y,x)=|x|^2-|y|^2+|x-y|^2$ gives
\begin{equation}\label{eq:discrete-energy-identity}
\begin{aligned}
 E^{n+1}-E^n
 &+\frac12\sum_{i=1}^2
       \bigl(|u_{n,i}-u_{n,i-1}|^2+|r_{n,i}-r_{n,i-1}|^2\bigr)\\
 &+\nu\tau_n\sum_{i=1}^2\sum_{j=1}^i
              a_{ij}(\nabla u_{n,j},\nabla u_{n,i})\\
 &+\gamma\tau_n\sum_{i=1}^2\sum_{j=1}^i
              a_{ij}r_{n,j}(r_{n,i}-1)
 =\tau_n\sum_{i=1}^2\sum_{j=0}^{i-1}
              \widehat a_{ij}\langle F_{n,j},u_{n,i}\rangle.
\end{aligned}
\end{equation}
The velocity quadratic form is bounded below by
$\lambda_I\sum_i\|u_{n,i}\|^2$. For the scalar term, write
$r_{n,j}=(r_{n,j}-1)+1$. Its row sums are $\eta$ and $1-\eta$,
so \cref{eq:incremental-coercivity} and Young's inequality imply
\[
 \sum_{i=1}^2\sum_{j=1}^i a_{ij}r_{n,j}(r_{n,i}-1)
 \ge\frac{\lambda_I}{2}\sum_{i=1}^2|r_{n,i}-1|^2
       -\frac{\eta^2+(1-\eta)^2}{2\lambda_I}.
\]
For each force row,
\[
 \left|\sum_{j=0}^{i-1}\widehat a_{ij}
                 \langle F_{n,j},u_{n,i}\rangle\right|
 \le\frac{\nu\lambda_I}{2}\|u_{n,i}\|^2
  +\frac{M_{F,V'}^2}{2\nu\lambda_I}
                  \left(\sum_{j=0}^{i-1}|\widehat a_{ij}|\right)^2.
\]
The two absolute row sums are $\eta$ and $2-\delta$.
Substitution into \cref{eq:discrete-energy-identity} proves the result.
\end{proof}

\begin{theorem}[Uniform $L^2$ bounds and absorption]\label{thm:discrete-l2}
Under \cref{eq:discrete-force-dual}, for every positive time grid,
\begin{equation}\label{eq:discrete-l2-bound}
 E^N\le Q_NE^0+\frac{C_0}{\kappa}(1-Q_N)\le M_E,
 \qquad Q_N=\prod_{n=0}^{N-1}(1+\kappa\tau_n)^{-1},\quad Q_0=1.
\end{equation}
For $0\le m<N$, the stages satisfy
\begin{equation}\label{eq:discrete-h1-window}
 \sum_{n=m}^{N-1}\tau_n\sum_{i=1}^2\|u_{n,i}\|^2
 \le\frac{2}{\nu\lambda_I}\bigl[M_E+C_0(t_N-t_m)\bigr].
\end{equation}
If $t_N\to\infty$, then $\limsup_{N\to\infty}E^N\le C_0/\kappa$.
In particular, every energy sublevel $E\le C_0/\kappa+\varepsilon$,
$\varepsilon>0$, absorbs all initial states of bounded energy.
If also $\tau_n\le\bar\tau<\infty$, there is a constant
$K_0=8M_E+4C_0\bar\tau$ such that
\begin{equation}\label{eq:stage-low-bounds}
 |u_{n,i}|^2\le K_0,\qquad |r_{n,i}-1|^2\le K_0,
 \qquad n\ge0,\quad i=0,1,2.
\end{equation}
\end{theorem}

\begin{proof}
Keep only the last-stage dissipation in
\cref{eq:discrete-energy-inequality} and use Poincare's inequality.
Since $u_{n,2}=u^{n+1}$ and $r_{n,2}=r^{n+1}$, this gives
\[
 (1+\kappa\tau_n)E^{n+1}\le E^n+C_0\tau_n.
\]
Subtracting $(1+\kappa\tau_n)C_0/\kappa$ and iterating proves
\cref{eq:discrete-l2-bound}. Summation of the full energy inequality
from $m$ to $N-1$, followed by $E^N\ge0$ and $E^m\le M_E$,
gives \cref{eq:discrete-h1-window}.

For absorption, the positive cross terms in the product yield
$Q_N\le(1+\kappa t_N)^{-1}$. Therefore
\[
 E^N\le\frac{C_0}{\kappa}
       +\frac{(E^0-C_0/\kappa)_+}{1+\kappa t_N},
 \qquad x_+=\max\{x,0\}.
\]
The second term tends to zero uniformly over bounded sets of initial
energies. No upper bound or ratio condition on the steps is needed
for this nodal conclusion.

Finally, if $\tau_n\le\bar\tau$, the increment terms in
\cref{eq:discrete-energy-inequality} imply
\[
 \sum_{i=1}^2\bigl(|u_{n,i}-u_{n,i-1}|^2
                   +|r_{n,i}-r_{n,i-1}|^2\bigr)
 \le2(M_E+C_0\bar\tau).
\]
Using $|x+y|^2\le2|x|^2+2|y|^2$ for the first-stage velocity and
shifted scalar proves \cref{eq:stage-low-bounds}. Stages $0$ and $2$
are already controlled by the nodal bound.
\end{proof}

By Poincare's inequality, \cref{eq:discrete-h1-window} also bounds
the full $H^1$ norm on every window of bounded physical length.
It is an integral-in-time estimate and does not alone provide a
step-size-independent pointwise $H^1$ bound. For a fixed step,
$Q_N=(1+\kappa\tau)^{-N}$, so the nodal transient decays exponentially
in physical time.

\subsection{Uniform \texorpdfstring{$H^1$}{H1} stability for a fixed time step}
\label{sec:discrete-h1}

We now assume
\begin{equation}\label{eq:h1-discrete-assumptions}
 u^0\in D(A),\qquad r^0=0,\qquad
 M_{F,H}:=\sup_{n\ge0,\ j=0,1}|F_{n,j}|<\infty,\qquad
 \tau_n=\tau\in(0,\bar\tau].
\end{equation}
Here $\bar\tau<\infty$ is fixed in advance. In this subsection, constants
$C$ may depend on $\Omega$, $\nu$, $\gamma$, $M_{F,H}$, $\|u^0\|$,
$|Au^0|$, and $\bar\tau$, but not on $n$, the particular $\tau$, or
the real-root branch. Constants $C_\theta$ may additionally depend on
a fixed positive parameter $\theta$. The preceding estimates apply
with $M_{F,V'}=\lambda_1^{-1/2}M_{F,H}$.

We first collect their low-order consequences. In addition to
\cref{eq:stage-low-bounds}, which bounds $|1-r_{n,i}^2|$, we have
\begin{align}
 \tau\bigl(\|u^n\|^2+\|u_{n,1}\|^2\bigr)&\le C,
       \label{eq:weighted-low-point}\\
 \tau\sum_{n=m}^{N-1}\bigl(\|u^n\|^2+\|u_{n,1}\|^2\bigr)
       &\le C\bigl(1+(N-m)\tau\bigr),\qquad 0\le m<N.
       \label{eq:weighted-low-window}
\end{align}
Indeed, the first-stage terms follow directly from
\cref{eq:discrete-h1-window}. For the old nodal terms with $m\ge1$,
use $u^n=u_{n-1,2}$ and shift the summation interval to
$m-1,\ldots,N-2$, whose length is unchanged for a fixed step.
When $m=0$, separate the initial term and use
$\tau\|u^0\|^2\le\bar\tau\|u^0\|^2$. Taking a single-step window
also proves \cref{eq:weighted-low-point}.

\begin{lemma}[Higher-order stage estimates]\label[lemma]{lem:h1-stage}
Under \cref{eq:h1-discrete-assumptions}, fix
$0<\theta<\nu\lambda_I/4$. Then
\begin{equation}\label{eq:h1-master}
\begin{aligned}
 \|u^{n+1}\|^2-\|u^n\|^2
 &+\nu\lambda_I\tau\bigl(|Au_{n,1}|^2+|Au^{n+1}|^2\bigr)\\
 &\le\theta\tau|Au^n|^2
       +C_\theta\tau\bigl(\|u^n\|^4+\|u_{n,1}\|^4+1\bigr),
\end{aligned}
\end{equation}
and the first stage satisfies
\begin{equation}\label{eq:h1-first-stage}
 \|u_{n,1}\|^2\le\|u^n\|^2+\frac\theta2\tau|Au^n|^2
                  +C_\theta\tau\|u^n\|^4+C_\theta\tau.
\end{equation}
\end{lemma}

\begin{proof}
By \cref{prop:stage-existence}, all stage velocities belong to $D(A)$.
We may therefore take the $H$ inner product of the velocity equation
with $2Au_{n,i}$. Summation over $i$ gives the gradient energy difference,
the nonnegative increments $\sum_i\|u_{n,i}-u_{n,i-1}\|^2$, and a
viscous term bounded below by
$2\nu\lambda_I\tau\sum_i|Au_{n,i}|^2$.

There are exactly three convective pairings, indexed by
$(i,j)=(1,0),(2,0),(2,1)$. Estimate \cref{eq:trilinear-a} and the
low-order bounds give, for each pair,
\begin{equation}\label{eq:h1-nonlinear-young}
\begin{aligned}
 2\bigl|(1-r_{n,i}^2)\widehat a_{ij}
                     (B_{n,j},Au_{n,i})\bigr|
 &\le C\|u_{n,j}\||Au_{n,j}|^{1/2}|Au_{n,i}|\\
 &\le\varepsilon|Au_{n,i}|^2+\sigma|Au_{n,j}|^2
                    +C_{\varepsilon,\sigma}\|u_{n,j}\|^4.
\end{aligned}
\end{equation}
Take $\varepsilon=\nu\lambda_I/4$ for all three terms,
$\sigma=\theta/2$ for the two terms with $j=0$, and
$\sigma=\nu\lambda_I/4$ for $(i,j)=(2,1)$.
Their sum uses at most
$(\nu\lambda_I/2)(|Au_{n,1}|^2+|Au^{n+1}|^2)$ of the current
viscous dissipation, leaving the old-node term $\theta|Au^n|^2$.
For the force, each row satisfies
\[
 2\left|\sum_{j<i}\widehat a_{ij}(F_{n,j},Au_{n,i})\right|
 \le\frac{\nu\lambda_I}{4}|Au_{n,i}|^2+C.
\]
After multiplication by $\tau$ and absorption, the coefficient of each
current $|Au|^2$ is at least $5\nu\lambda_I/4$.
Discarding the increments and reducing this coefficient to
$\nu\lambda_I$ proves \cref{eq:h1-master}.
For the first stage alone, the same estimates leave the positive
coefficient $2\nu\eta-\nu\lambda_I/2$ on $\tau|Au_{n,1}|^2$.
Dropping that dissipation and the increment proves
\cref{eq:h1-first-stage}.
\end{proof}

\begin{theorem}[Uniform $H^1$ stability]\label{thm:discrete-h1}
Under \cref{eq:h1-discrete-assumptions}, every solution sequence of
\cref{eq:sdirk-mr-scheme} satisfies
\begin{equation}\label{eq:discrete-h1-uniform}
 \sup_{n\ge0,\ i=0,1,2}\|u_{n,i}\|^2\le C.
\end{equation}
Moreover, for every $0\le m<N$,
\begin{equation}\label{eq:discrete-h2-window}
 \tau\sum_{n=m}^{N-1}\sum_{i=1}^2|Au_{n,i}|^2
 \le C\bigl(1+(N-m)\tau\bigr).
\end{equation}
The constants have the dependencies specified above and are uniform
for $0<\tau\le\bar\tau$. No additional small-step condition is required.
\end{theorem}

\begin{proof}
The old-node term in \cref{eq:h1-master} prevents a direct estimate
of the gradient energy alone. We include it in a modified energy,
prove a window bound for that energy without using the desired
pointwise estimate, and then average an explicit recurrence over a
fixed physical-time window.

\emph{Modified energy and its recurrence.}
Fix $\theta=\nu\lambda_I/8$ and define
\begin{equation}\label{eq:modified-h1-energy}
 X_n=\|u^n\|^2+\theta\tau|Au^n|^2.
\end{equation}
Rearranging \cref{eq:h1-master} gives
\begin{equation}\label{eq:modified-h1-dissipation}
\begin{aligned}
 X_{n+1}-X_n
 &+(\nu\lambda_I-\theta)\tau
                   \bigl(|Au_{n,1}|^2+|Au^{n+1}|^2\bigr)\\
 &\le C\tau\bigl(\|u^n\|^4+\|u_{n,1}\|^4+1\bigr).
\end{aligned}
\end{equation}
By \cref{eq:weighted-low-point},
$\tau\|u^n\|^4\le C\|u^n\|^2\le CX_n$.
Thus \cref{eq:h1-first-stage} implies
\begin{equation}\label{eq:internal-h1-control}
 \|u_{n,1}\|^2\le CX_n+C\tau.
\end{equation}
Multiplication by $\|u_{n,1}\|^2$ and another use of
\cref{eq:weighted-low-point} yield
$\|u_{n,1}\|^4\le C\|u_{n,1}\|^2X_n+C$.
Consequently, \cref{eq:modified-h1-dissipation} implies
\begin{equation}\label{eq:explicit-h1-recurrence}
 X_{n+1}\le
 \bigl[1+C\tau(\|u^n\|^2+\|u_{n,1}\|^2)\bigr]X_n+C\tau.
\end{equation}
Choose a common constant for these estimates and hold it fixed during
iteration. The coefficient multiplies $X_n$, so no denominator
positivity condition involving the time step is needed.

\emph{A window bound for the modified energy.}
The gradient part of $\tau\sum_n X_n$ is already bounded by
\cref{eq:weighted-low-window}. The remaining term contains
$\tau^2|Au^n|^2$, which we must control first. Set
$W_n=\tau^2|Au^n|^2$. Multiplying \cref{eq:h1-master} by $\tau$,
discarding the new gradient term and first-stage dissipation, and
dividing by $\nu\lambda_I$ gives
\begin{equation}\label{eq:weighted-stokes-recurrence}
 W_{n+1}\le\tfrac18 W_n
    +C\tau\|u^n\|^2
    +C\tau^2\bigl(\|u^n\|^4+\|u_{n,1}\|^4+1\bigr).
\end{equation}
The last two terms are uniformly bounded by
\cref{eq:weighted-low-point} and $\tau\le\bar\tau$.
A geometric-series argument, with
$W_0\le\bar\tau^2|Au^0|^2$, therefore gives $\sup_nW_n\le C$.
For the window sum, retain the time weights and observe that
\[
 \tau^2\|u^n\|^4\le C\tau\|u^n\|^2,\qquad
 \tau^2\|u_{n,1}\|^4\le C\tau\|u_{n,1}\|^2.
\]
Using \cref{eq:weighted-low-window} and
$(N-m)\tau^2\le\bar\tau(N-m)\tau$, the sum of the driving terms
in \cref{eq:weighted-stokes-recurrence} is bounded by
$C(1+(N-m)\tau)$. Since
$\sum_{n=m}^{N-1}W_{n+1}=\sum_{n=m}^{N-1}W_n-W_m+W_N$, we obtain
\[
 \tfrac78\sum_{n=m}^{N-1}W_n
 \le W_m-W_N+C\bigl(1+(N-m)\tau\bigr)
 \le C\bigl(1+(N-m)\tau\bigr).
\]
It follows that
\begin{equation}\label{eq:modified-h1-window}
 \tau\sum_{n=m}^{N-1}X_n
 =\tau\sum_{n=m}^{N-1}\|u^n\|^2
       +\theta\sum_{n=m}^{N-1}W_n
 \le C\bigl(1+(N-m)\tau\bigr).
\end{equation}
This argument uses neither a pointwise $H^1$ bound nor an $H^2$
time-window bound.

\emph{Uniform Gronwall estimate and the initial interval.}
Iteration of \cref{eq:explicit-h1-recurrence} from any $k<N$, using
$1+x\le e^x$ and \cref{eq:weighted-low-window}, yields
\begin{equation}\label{eq:h1-iteration}
 X_N\le e^{C(1+(N-k)\tau)}\bigl[X_k+C(N-k)\tau\bigr].
\end{equation}
Fix $T_w>0$ independently of $\tau$. If $N\tau\ge T_w$, set
$m=N-\lceil T_w/\tau\rceil$ and $L=(N-m)\tau$, so that
$T_w\le L<T_w+\bar\tau$. For each $k=m,\ldots,N-1$, the preceding
estimate is bounded by $e^{C(1+L)}(X_k+CL)$.
Multiply by $\tau$, sum over $k$, and divide by $L$. By
\cref{eq:modified-h1-window},
\[
 X_N\le e^{C(1+L)}
        \left[\frac{\tau}{L}\sum_{k=m}^{N-1}X_k+CL\right]
 \le e^{C(1+T_w+\bar\tau)}
        \left[\frac{C(1+T_w+\bar\tau)}{T_w}
                         +C(T_w+\bar\tau)\right].
\]
For $N\tau<T_w$, apply \cref{eq:h1-iteration} with $k=0$ and use
\[
 X_0\le\|u^0\|^2+\theta\bar\tau|Au^0|^2.
\]
The case $N=0$ follows directly from the same initial bound.
Fixing $T_w$ once and for all gives $\sup_nX_n\le C$.
The nodal gradient norm is bounded by $X_n$, and
\cref{eq:internal-h1-control} bounds the first stage.
The second stage is the next node. This proves
\cref{eq:discrete-h1-uniform}.

\emph{Higher-order dissipation on time windows.}
We may now use the proved pointwise bound to write
$\|u^n\|^4+\|u_{n,1}\|^4
 \le C(\|u^n\|^2+\|u_{n,1}\|^2)$.
Summing \cref{eq:modified-h1-dissipation} from $m$ to $N-1$ gives
\[
\begin{aligned}
 (\nu\lambda_I-\theta)\tau
      \sum_{n=m}^{N-1}\bigl(|Au_{n,1}|^2+|Au^{n+1}|^2\bigr)
 &\le X_m-X_N\\
 &\quad+C\tau\sum_{n=m}^{N-1}
                 (\|u^n\|^2+\|u_{n,1}\|^2+1)\\
 &\le C\bigl(1+(N-m)\tau\bigr).
\end{aligned}
\]
Because $\nu\lambda_I-\theta>0$, this proves
\cref{eq:discrete-h2-window}.
\end{proof}

By \cref{eq:stokes-estimates}, the last theorem gives a full $H^1$
uniform bound and an $H^2$ time-window bound. Its constants retain
initial-data dependence, so it does not assert an $H^1$ absorbing
radius independent of the initial data. Moreover, $X_n\le C$ only
yields $|Au^n|^2\le C/(\theta\tau)$; a step-size-independent pointwise
$H^2$ estimate does not follow. The assumption $u^0\in D(A)$ is used
both to justify the stage tests and to control the initial weighted
Stokes energy. The fixed-step hypothesis is used when shifting window
sums and forming $X_n$ and $W_n$; the $H^1$ conclusion here is not
claimed for arbitrary variable steps.

\section{Refined auxiliary-variable estimates and finite-time convergence}
\label{sec:convergence}

In this section, we first prove that the auxiliary variable remains
of order $\tau$ uniformly in time. Its contribution to the velocity
coefficient is therefore of order $\tau^2$. We then construct reference
stages and establish second-order convergence of the velocity on a
fixed finite time interval. Throughout this section the time step is
constant, $0<\tau\le\bar\tau$, and the assumptions in
\cref{eq:h1-discrete-assumptions} are in force. The estimates apply to
all real-root branches of the scheme.

\subsection{A refined estimate for the auxiliary variable}
\label{sec:auxiliary-refined}

\begin{theorem}[Uniform first-order auxiliary-variable bound]
\label{thm:auxiliary-refined}
Under \cref{eq:h1-discrete-assumptions}, there is a constant $C$,
with the same allowed dependencies as in \cref{thm:discrete-h1}, such that
\begin{equation}\label{eq:auxiliary-refined}
 \sup_{n\ge0,\ i=0,1,2}|r_{n,i}|\le C\tau,
 \qquad
 \sup_{n\ge0,\ i=0,1,2}|G_u(r_{n,i})-1|\le C\tau^2.
\end{equation}
In particular, for $0\le j<i\le2$,
\begin{equation}\label{eq:auxiliary-velocity-perturbation}
 \|r_{n,i}^2B(u_{n,j},u_{n,j})\|_{V'}\le C\tau^2.
\end{equation}
The constants are independent of the final time, the particular step
size, and the selected roots.
\end{theorem}

\begin{proof}
We use the established $H^1$ bound and $H^2$ window estimate to control
the accumulated scalar forcing. Exponential damping of the scalar
recurrence then prevents this forcing from accumulating over arbitrarily
long times.

\emph{Window estimates for convection and stage differences.}
By \cref{eq:trilinear-a} and \cref{thm:discrete-h1},
\[
 |B_{n,j}|^2\le C|u_{n,j}|\|u_{n,j}\|^2|Au_{n,j}|
              \le C(1+|Au_{n,j}|^2),\qquad j=0,1.
\]
The $H^2$ window bound also controls the old nodes: for $m\ge1$,
shift $u^n=u_{n-1,2}$; for $m=0$, separate the term
$\tau|Au^0|^2\le\bar\tau|Au^0|^2$.
It follows that, for every $0\le m<N$,
\begin{equation}\label{eq:convection-window}
 \tau\sum_{n=m}^{N-1}\sum_{j=0}^1|B_{n,j}|^2
 \le C\bigl(1+(N-m)\tau\bigr).
\end{equation}
Dividing the velocity stages by $\tau$ and using the bounded
coefficients $|1-r_{n,i}^2|$ from \cref{eq:stage-low-bounds} gives
\[
 \sum_{i=1}^2\left|\frac{u_{n,i}-u_{n,i-1}}{\tau}\right|^2
 \le C\left(\sum_{i=1}^2|Au_{n,i}|^2
              +\sum_{j=0}^1|B_{n,j}|^2+1\right).
\]
Hence
\begin{equation}\label{eq:stage-difference-window}
 \tau\sum_{n=m}^{N-1}\sum_{i=1}^2
       \left|\frac{u_{n,i}-u_{n,i-1}}{\tau}\right|^2
 \le C\bigl(1+(N-m)\tau\bigr).
\end{equation}
These are window estimates, not pointwise bounds on the stage
difference quotients.

\emph{The accumulated scalar defect.}
Define the scalar work defect and its contribution to the scalar equation by
\[
 \mathcal D_{n,i}=\sum_{j=0}^{i-1}\widehat a_{ij}(B_{n,j},u_{n,i}),
 \qquad h_{n,i}=(1+r_{n,i})\mathcal D_{n,i}.
\]
All inner products are in $H$, since the stages belong to $D(A)$.
Skew-symmetry gives $(B_{n,j},u_{n,j})=0$, and thus
\[
 \mathcal D_{n,i}=\sum_{j=0}^{i-1}\widehat a_{ij}
                         (B_{n,j},u_{n,i}-u_{n,j}).
\]
Expressing each difference as a sum of adjacent stage differences yields
\[
 \sum_{i=1}^2|\mathcal D_{n,i}|
 \le C\tau\left(\sum_{j=0}^1|B_{n,j}|^2\right)^{1/2}
       \left(\sum_{i=1}^2
          \left|\frac{u_{n,i}-u_{n,i-1}}{\tau}\right|^2\right)^{1/2}.
\]
Multiply by $\tau$ and sum over a window. Weighted Cauchy--Schwarz,
\cref{eq:convection-window,eq:stage-difference-window}, and the existing
bound on $|1+r_{n,i}|$ give
\begin{equation}\label{eq:scalar-driving-window}
 \tau\sum_{n=m}^{N-1}\bigl(|h_{n,1}|+|h_{n,2}|\bigr)
 \le C_h\tau\bigl(1+(N-m)\tau\bigr).
\end{equation}
Fix this constant $C_h$ for the remainder of the proof. This estimate
uses only the previously proved $O(1)$ bound on the scalar, not the
first-order bound being proved.

\emph{Damped propagation of the scalar forcing.}
Set
\[
 s(z)=\frac1{1+\eta z},\qquad
 \mathcal R(z)=\frac{1-(1-2\eta)z}{(1+\eta z)^2}.
\]
Eliminating the first scalar stage gives the exact identities
\begin{equation}\label{eq:scalar-damped-recurrence}
\begin{aligned}
 r_{n,1}&=s(\gamma\tau)r^n-\tau s(\gamma\tau)h_{n,1},\\
 r^{n+1}&=\mathcal R(\gamma\tau)r^n
       -\tau\mathcal R(\gamma\tau)h_{n,1}
       -\tau s(\gamma\tau)h_{n,2}.
\end{aligned}
\end{equation}
Here $h_{n,i}$ is defined from the computed solution, so this identity
does not assume that the nonlinear stages are independent linear problems.
Clearly $0<s(z)\le1$ for $z\ge0$. If the numerator of $\mathcal R(z)$
is nonnegative, then
\[
 1-|\mathcal R(z)|=\frac{z+\eta^2z^2}{(1+\eta z)^2}.
\]
If it is negative, then
\[
 1-|\mathcal R(z)|=\frac{2+(4\eta-1)z+\eta^2z^2}{(1+\eta z)^2}.
\]
Since $0<4\eta-1<1$, both cases imply
\begin{equation}\label{eq:scalar-uniform-decay}
 |\mathcal R(\gamma\tau)|\le1-c\tau\le e^{-c\tau},
 \qquad c=\frac{(4\eta-1)\gamma}{(1+\eta\gamma\bar\tau)^2}>0.
\end{equation}
The first inequality also ensures $1-c\tau\ge0$ on the prescribed
step range. Using $r^0=0$ in \cref{eq:scalar-damped-recurrence},
\[
 |r^N|\le\tau\sum_{n=0}^{N-1}e^{-c(N-1-n)\tau}
                           \bigl(|h_{n,1}|+|h_{n,2}|\bigr).
\]
Fix a physical window length $T_w>0$, and set
$K=\lceil T_w/\tau\rceil$, so $T_w\le K\tau<T_w+\bar\tau$.
Partition this sum backwards from $n=N-1$ into blocks of at most $K$
steps. On the $\ell$th block, counting from zero, the exponential is
at most $e^{-c\ell K\tau}$, while \cref{eq:scalar-driving-window}
bounds its unweighted contribution by $C_h\tau(1+K\tau)$.
The earliest, possibly incomplete, block satisfies the same bound.
Therefore
\[
 |r^N|\le C_h\tau(1+T_w+\bar\tau)
                   \sum_{\ell=0}^{\infty}e^{-c\ell T_w}
       =\frac{C_h(1+T_w+\bar\tau)}{1-e^{-cT_w}}\tau.
\]
This also holds for $N=0$. The first identity in
\cref{eq:scalar-damped-recurrence}, together with the single-step case
of \cref{eq:scalar-driving-window}, gives
$|r_{n,1}|\le|r^n|+\tau|h_{n,1}|\le C\tau$.
Stages $0$ and $2$ are nodes. Finally,
$|G_u(r)-1|=r^2$ proves \cref{eq:auxiliary-refined}, and
$\|B(v,v)\|_{V'}\le C\|v\|^2$ proves
\cref{eq:auxiliary-velocity-perturbation}.
\end{proof}

\subsection{Reference stages and local consistency}
\label{sec:reference-consistency}

Fix $T>0$. For the consistency analysis we impose the sufficient
regularity assumptions
\begin{equation}\label{eq:exact-solution-regularity}
\begin{aligned}
 u&\in W^{3,\infty}(0,T;H)\cap W^{2,\infty}(0,T;D(A))
                  \cap W^{1,\infty}(0,T;D(A^2)),\\
 F&\in W^{2,\infty}(0,T;H).
\end{aligned}
\end{equation}
Here $D(A^2)=\{v\in D(A):Av\in D(A)\}$, endowed with its graph
norm. In particular, we use
\[
 \|v\|_{D(A)}^2=|v|^2+|Av|^2,\qquad
 \|v\|_{D(A^2)}^2=|v|^2+|Av|^2+|A^2v|^2.
\]
These assumptions include Stokes boundary compatibility and are not
identified with ordinary spatial $H^4$ regularity. They concern the
exact solution, not the numerical stages. On every complete time step
$t_{n+1}\le T$, take the force samples
\begin{equation}\label{eq:convergence-force-sampling}
 F_{n,0}=F(t_n),\qquad F_{n,1}=F(t_n+\eta\tau).
\end{equation}
Constants $C_T$ may depend on the norms in
\cref{eq:exact-solution-regularity}, $T$, and the fixed data, but not on
$n$ or the particular $\tau$.

Define $q(t)=F(t)-B(u(t),u(t))$, so that $u_t=-\nu Au+q$.
The reference stages are
\begin{equation}\label{eq:reference-stages}
\begin{aligned}
 U_{n,0}&=u(t_n),\qquad U_{n,2}=u(t_{n+1}),\\
 (I+\eta\nu\tau A)U_{n,1}&=u(t_n)+\eta\tau q(t_n).
\end{aligned}
\end{equation}
Thus the first reference stage satisfies the uncorrected velocity
stage exactly, while the second is the exact endpoint. Define the
second-stage increment residual by
\begin{equation}\label{eq:reference-residual}
\begin{aligned}
 \mathcal R_n
 &=U_{n,2}-U_{n,1}
       +\nu\tau\bigl[(1-2\eta)AU_{n,1}+\eta AU_{n,2}\bigr]\\
 &\quad-\tau\sum_{j=0}^1\widehat a_{2j}
                    \bigl[F_{n,j}-B(U_{n,j},U_{n,j})\bigr].
\end{aligned}
\end{equation}
This residual appears with a positive sign on the right-hand side of
the reference second-stage equation.

\begin{lemma}[Reference-stage consistency]\label[lemma]{lem:reference-consistency}
Under \cref{eq:exact-solution-regularity,eq:convergence-force-sampling},
for every $t_{n+1}\le T$,
\begin{equation}\label{eq:reference-consistency-bounds}
 \|U_{n,1}-u(t_n+\eta\tau)\|_{D(A)}\le C_T\tau^2,
 \qquad |\mathcal R_n|\le C_T\tau^3.
\end{equation}
Moreover, $\max_{i=0,1,2}\|U_{n,i}\|_{D(A)}\le C_T$.
\end{lemma}

\begin{proof}
Set $S_\tau=(I+\eta\nu\tau A)^{-1}$. The spectral representation of
the positive self-adjoint Stokes operator shows that $S_\tau$ commutes
with $A$ on its domain and is a contraction in $H$, $D(A)$, and
$D(A^2)$. In this proof, $O_H(\tau^k)$ denotes a remainder whose
$H$ norm is at most $C_T\tau^k$.

The equation for $u$ gives the exact identity
\begin{equation}\label{eq:reference-resolvent-identity}
 U_{n,1}=u(t_n)+\eta\tau S_\tau u_t(t_n).
\end{equation}
The nodal values $u_t(t_n)$ belong to $D(A^2)$ with a common bound:
this follows from $H$ continuity of $u_t$, its essential bound in
$D(A^2)$, and weak compactness in that graph space.
Using $(S_\tau-I)v=-\eta\nu\tau AS_\tau v$ and Taylor's formula
in $D(A)$, we obtain
\[
\begin{aligned}
 \|U_{n,1}-u(t_n+\eta\tau)\|_{D(A)}
 &\le\eta^2\tau^2\left(
       \nu\|u_t\|_{L^\infty(0,T;D(A^2))}
       +\tfrac12\|u_{tt}\|_{L^\infty(0,T;D(A))}\right).
\end{aligned}
\]
The uniform graph-norm bound on $U_{n,1}$ follows directly from
\cref{eq:reference-resolvent-identity}.

To estimate the endpoint residual, expand the same identity once more
in $H$. The contraction bounds give
\begin{align*}
 U_{n,1}&=u(t_n)+\eta\tau u_t(t_n)
                 -\eta^2\nu\tau^2 Au_t(t_n)+O_H(\tau^3),\\
 AU_{n,1}&=Au(t_n)+\eta\tau Au_t(t_n)+O_H(\tau^2).
\end{align*}
For the second line we apply $A$ to
\cref{eq:reference-resolvent-identity} directly; its remainder is
$-\eta^2\nu\tau^2A^2S_\tau u_t(t_n)$. No derivative is applied
to a remainder controlled only in $H$.
The estimate
$|B(v,w)|\le C\|v\|_{D(A)}\|w\|_{D(A)}$ implies
$q\in W^{2,\infty}(0,T;H)$ and local Lipschitz continuity of
$B(v,v)$ from $D(A)$ to $H$. Hence
\[
 F_{n,1}-B(U_{n,1},U_{n,1})
     =q(t_n)+\eta\tau q'(t_n)+O_H(\tau^2).
\]
Taylor expansion of the exact endpoint, using
$u_{tt}=-\nu Au_t+q'$, gives
\begin{equation}\label{eq:exact-endpoint-expansion}
\begin{aligned}
 (I+\eta\nu\tau A)u(t_n+\tau)
 &=u(t_n)+\tau[-\nu(1-\eta)Au(t_n)+q(t_n)]\\
 &\quad+\tau^2[\nu(\eta-\tfrac12)Au_t(t_n)
                         +\tfrac12q'(t_n)]+O_H(\tau^3).
\end{aligned}
\end{equation}
The third-order remainder is controlled by $u_{ttt}$ in $H$ and
$Au_{tt}$ in $H$. On the other hand, the reference right-hand side is
\begin{equation}\label{eq:reference-endpoint-expansion}
\begin{aligned}
 &U_{n,1}-\nu\tau(1-2\eta)AU_{n,1}
       +\tau\sum_{j=0}^1\widehat a_{2j}
                 [F_{n,j}-B(U_{n,j},U_{n,j})]\\
 &=u(t_n)+\tau[-\nu(1-\eta)Au(t_n)+q(t_n)]\\
 &\quad+\tau^2[-\nu\eta(1-\eta)Au_t(t_n)
                         +\eta(1-\delta)q'(t_n)]+O_H(\tau^3).
\end{aligned}
\end{equation}
The coefficient identities
\[
 \eta(1-\eta)=\tfrac12-\eta,\qquad
 \eta(1-\delta)=\tfrac12
\]
make the second-order terms agree. Subtracting
\cref{eq:reference-endpoint-expansion} from
\cref{eq:exact-endpoint-expansion} proves $|\mathcal R_n|\le C_T\tau^3$.
\end{proof}

\subsection{Finite-time velocity error estimates}
\label{sec:velocity-convergence}

\begin{theorem}[Second-order velocity convergence]\label{thm:velocity-convergence}
Assume \cref{eq:h1-discrete-assumptions,eq:exact-solution-regularity,eq:convergence-force-sampling}, with $u^0=u(0)$ and $r^0=0$.
Then, for $N_T=\lfloor T/\tau\rfloor$,
\begin{equation}\label{eq:velocity-convergence}
\begin{aligned}
 \max_{0\le n\le N_T}|u^n-u(t_n)|&\le C_T\tau^2,\\
 \left(\tau\sum_{n=0}^{N_T-1}\sum_{i=1}^2
       \|u_{n,i}-u(t_n+c_i\tau)\|^2\right)^{1/2}
       &\le C_T\tau^2.
\end{aligned}
\end{equation}
The constant is independent of $\tau\in(0,\bar\tau]$ and of the
real-root branch. It may depend on $T$ and the specified exact-solution
norms, in addition to the data entering the stability estimates.
\end{theorem}

\begin{proof}
Compare the numerical stages with \cref{eq:reference-stages} and set
\[
 e_{n,i}=u_{n,i}-U_{n,i},\qquad
 e^n=e_{n,0}=u^n-u(t_n),\qquad e^{n+1}=e_{n,2}.
\]
For $0\le j<i\le2$, define the full nonlinear difference
\[
 \mathcal N_{n,i,j}=(1-r_{n,i}^2)B(u_{n,j},u_{n,j})
                                  -B(U_{n,j},U_{n,j}).
\]
Subtracting the reference equations from the numerical equations yields
\begin{equation}\label{eq:velocity-error-stages}
\begin{aligned}
 e_{n,1}-e^n+\nu\eta\tau Ae_{n,1}
      &=-\eta\tau\mathcal N_{n,1,0},\\
 e_{n,2}-e_{n,1}
      +\nu\tau[(1-2\eta)Ae_{n,1}+\eta Ae_{n,2}]
      &=-\tau\sum_{j=0}^1\widehat a_{2j}\mathcal N_{n,2,j}
         -\mathcal R_n.
\end{aligned}
\end{equation}
The force samples cancel, and the sign of the residual follows from
\cref{eq:reference-residual}.

\emph{Nonlinear error estimates.}
The velocity part of $\mathcal N_{n,i,j}$ decomposes as
\[
 B(u_{n,j},u_{n,j})-B(U_{n,j},U_{n,j})
       =B(e_{n,j},u_{n,j})+B(U_{n,j},e_{n,j}).
\]
Move the derivative onto the test error by skew-symmetry. The uniform
$V$ bounds on numerical and reference velocities and the two-dimensional
$L^4$ interpolation inequality then give
\[
\begin{aligned}
 &\left|\langle B(u_{n,j},u_{n,j})-B(U_{n,j},U_{n,j}),e_{n,i}\rangle\right|\\
 &\qquad\le C_T|e_{n,j}|^{1/2}\|e_{n,j}\|^{1/2}\|e_{n,i}\|.
\end{aligned}
\]
Two applications of Young's inequality control the current and old
stage gradients with independently chosen positive coefficients.
The remaining term satisfies, by
\cref{eq:auxiliary-velocity-perturbation},
$|\langle r_{n,i}^2B_{n,j},e_{n,i}\rangle|
 \le C\tau^2\|e_{n,i}\|$.
Consequently, for arbitrary $\varepsilon_1,\varepsilon_2>0$,
\begin{equation}\label{eq:nonlinear-error-bound}
 |\langle\mathcal N_{n,i,j},e_{n,i}\rangle|
 \le\varepsilon_1\|e_{n,i}\|^2+\varepsilon_2\|e_{n,j}\|^2
       +C_{T,\varepsilon_1,\varepsilon_2}|e_{n,j}|^2
       +C_{T,\varepsilon_1}\tau^4.
\end{equation}
The old-stage gradient term is retained; it cannot be replaced by an
$L^2$ error bound using only the available $H^1$ velocity bounds.

\emph{Stage energy estimates.}
Test the two equations in \cref{eq:velocity-error-stages} with
$2e_{n,1}$ and $2e_{n,2}$, respectively, and add them.
Polarization and \cref{eq:incremental-coercivity} give
$|e^{n+1}|^2-|e^n|^2$, nonnegative increment terms, and at least
$2\nu\lambda_I\tau(\|e_{n,1}\|^2+\|e^{n+1}\|^2)$ of dissipation.
For any prescribed $b>0$, use \cref{eq:nonlinear-error-bound} with
\[
 \varepsilon_1=\frac{\nu\lambda_I}{4(2-\delta)},\qquad
 \varepsilon_2=
 \begin{cases}
 b/[2(1+\eta)],&j=0,\\
 \nu\lambda_I/[4(1-\delta)],&(i,j)=(2,1).
 \end{cases}
\]
After multiplying by the actual weights $2\eta$, $2$, and
$2(1-\delta)$, the old-node gradient coefficient is $b$.
The first-stage gradient coefficient is at most $\nu\lambda_I$,
and the second-stage coefficient is $\nu\lambda_I/2$.
For the residual, Poincare's inequality and
\cref{eq:reference-consistency-bounds} give
\[
 2|(\mathcal R_n,e_{n,2})|
 \le\frac{\nu\lambda_I}{2}\tau\|e^{n+1}\|^2
       +\frac{2|\mathcal R_n|^2}{\nu\lambda_I\lambda_1\tau}
 \le\frac{\nu\lambda_I}{2}\tau\|e^{n+1}\|^2+C_T\tau^5.
\]
Absorbing these terms and discarding the increments yields
\begin{equation}\label{eq:velocity-error-energy}
\begin{aligned}
 |e^{n+1}|^2-|e^n|^2
 &+\nu\lambda_I\tau(\|e_{n,1}\|^2+\|e^{n+1}\|^2)\\
 &\le b\tau\|e^n\|^2
       +C_{T,b}\tau(|e^n|^2+|e_{n,1}|^2)+C_{T,b}\tau^5.
\end{aligned}
\end{equation}
Testing the first stage alone and taking
$\varepsilon_1=\nu/2$, $\varepsilon_2=b_1/(2\eta)$ in
\cref{eq:nonlinear-error-bound} gives, for any $b_1>0$,
\begin{equation}\label{eq:first-stage-error}
 |e_{n,1}|^2\le(1+C_{T,b_1}\tau)|e^n|^2
                 +b_1\tau\|e^n\|^2+C_{T,b_1}\tau^5.
\end{equation}
Here a positive first-stage dissipation remains before it is discarded;
no current-stage $L^2$ term needs to be moved to the left.

\emph{A closed error-energy recurrence.}
Fix $\vartheta=\nu\lambda_I/4$, take $b=\vartheta/2$ and
$b_1=\vartheta$, and introduce
\[
 \mathcal E_n=|e^n|^2+\vartheta\tau\|e^n\|^2.
\]
Choose a common constant $C_T$ for
\cref{eq:velocity-error-energy,eq:first-stage-error} at these fixed
parameters. Then
$|e_{n,1}|^2\le(1+C_T\tau)\mathcal E_n+C_T\tau^5$.
Rewriting \cref{eq:velocity-error-energy} in terms of
$\mathcal E_n$ leaves $(b-\vartheta)\tau\|e^n\|^2\le0$ on the
right, while the new-node dissipation has coefficient
$\nu\lambda_I-\vartheta=3\nu\lambda_I/4$.
Substitution of the first-stage estimate, followed by
$\tau\le\bar\tau$, therefore gives a fixed constant $K_T$ such that
\begin{equation}\label{eq:closed-error-recurrence}
 \mathcal E_{n+1}
   +\frac{\nu\lambda_I}{2}\tau
                         (\|e_{n,1}\|^2+\|e^{n+1}\|^2)
 \le(1+K_T\tau)\mathcal E_n+K_T\tau^5.
\end{equation}
For example, $K_T=1+2C_T+C_T^2\bar\tau$ suffices.
This is an explicit recurrence in the old error energy and introduces
no additional small-step restriction.

\emph{Finite-time iteration and stage-error recovery.}
Exact initialization gives $\mathcal E_0=0$. For $N\tau\le T$,
iteration of \cref{eq:closed-error-recurrence} without its dissipation
implies
\[
 \mathcal E_N\le K_T\tau^5\sum_{k=0}^{N-1}(1+K_T\tau)^k
       \le K_TT e^{K_TT}\tau^4.
\]
Keeping the dissipation and summing the same recurrence gives
\[
 \frac{\nu\lambda_I}{2}\tau
       \sum_{n=0}^{N-1}\sum_{i=1}^2\|e_{n,i}\|^2
 \le K_T\tau\sum_{n=0}^{N-1}\mathcal E_n+K_TN\tau^5
 \le C_T\tau^4.
\]
Since $U_{n,2}=u(t_{n+1})$, the nodal part of
\cref{eq:velocity-convergence} follows immediately.
For the first stage,
\[
 u_{n,1}-u(t_n+\eta\tau)
       =e_{n,1}+[U_{n,1}-u(t_n+\eta\tau)].
\]
The graph-norm estimate in \cref{eq:reference-consistency-bounds}
controls the bracket in $V$ by $C_T\tau^2$. Thus
\[
 \tau\sum_{n=0}^{N-1}\sum_{i=1}^2
       \|u_{n,i}-u(t_n+c_i\tau)\|^2
 \le2\tau\sum_{n=0}^{N-1}\sum_{i=1}^2\|e_{n,i}\|^2
       +C_T(N\tau)\tau^4
 \le C_T\tau^4.
\]
Take $N=N_T$ and square roots to finish the proof. If $N_T=0$,
the nodal error is zero and the stage sum is empty.
\end{proof}

The auxiliary-variable bound is uniform for all time, whereas the
velocity convergence constant may grow with $T$. The stage estimate
in \cref{eq:velocity-convergence} is a time-weighted $H^1$ estimate;
it does not assert pointwise second-order $H^1$ convergence. The proof
uses the numerical $H^1$ bound and $H^2$ window estimate, together with
the separate exact-solution regularity in
\cref{eq:exact-solution-regularity}. It does not require a
step-size-independent pointwise $H^2$ bound on the numerical velocity.
Pressure error estimates and variable-step convergence are not included
in this theorem.

\section{Conclusions}
\label{sec:conclusions}

Some conclusions here.


\appendix
\section{An example appendix} 
\lipsum[71]

\begin{lemma}
Test Lemma.
\end{lemma}


\section*{Acknowledgments}
We would like to acknowledge the assistance of volunteers in putting
together this example manuscript and supplement.

\bibliographystyle{siamplain}
\bibliography{references}
\end{document}
